\documentclass[12pt]{ctexart}
\usepackage[a4paper,margin=2.5cm]{geometry}
\usepackage{graphicx}
\usepackage{longtable}
\usepackage{booktabs}
\usepackage{fancyhdr}
\newcommand{\SoftwareName}{软件名称待填写}
\newcommand{\SoftwareVersion}{V1.0}
\pagestyle{fancy}
\fancyhf{}
\lhead{\SoftwareName}
\rhead{\SoftwareVersion}
\cfoot{\thepage}
\title{\SoftwareName{} 软件说明书}
\author{}
\date{\SoftwareVersion}
\begin{document}
\maketitle
\section{软件概述}
说明软件用途、目标用户和核心价值。
\section{用途和适用范围}
说明软件适用场景、边界和运行前提。
\section{运行环境}
说明硬件环境、操作系统、依赖环境和部署方式。
\section{系统架构}
插入模块架构图、部署图和数据流说明。
\section{功能模块}
逐项说明功能模块、输入、处理和输出。
\section{数据输入输出}
说明主要数据项、格式、来源和输出结果。
\section{用户角色与权限}
说明用户类型、权限范围和操作边界。
\section{操作步骤}
逐步说明主要操作流程，并插入真实截图或标注为界面示例的原型截图。
\section{异常处理}
说明错误提示、恢复方式和数据保护措施。
\section{部署和维护}
说明安装、配置、备份、升级和维护注意事项。
\section{版本说明}
说明当前申请版本对应的功能范围。
\section{提交检查清单}
列出申请表素材、软件说明书、源程序、证明文件和视觉材料的准备状态。
\end{document}
